% !TeX root = ../main.tex
\section{问题重述}
\subsection{问题背景}
随着分布式光伏和储能技术的发展，微网逐渐成为组织区域内部发电、储电和用电的重要形式。微网通常将分布式电源、用户负荷、储能设备、能量转换装置及监控保护系统整合为一个可协调控制的小型发配电系统\cite{ref1}。
在智能小区中，微网能够进一步连接小区光伏、居民负荷和储能设备，并通过联络线路与外部电网进行电能交换，从而实现本地能源的统一管理\cite{ref2}。
这种运行方式既有利于提高分布式能源的就地利用水平，也能够增强小区供电的灵活性，为后续制定外网购电与储能充放电策略提供基础。

然而，小区微网的运行环境具有明显的不确定性。光伏发电容易受到天气和光照条件影响，居民用电需求也会随用户行为发生变化，使实际发电量和负荷可能偏离预测结果\cite{ref3}。
与此同时，外部电网电价受到能源供需、市场环境及政策变化等因素影响，其波动会改变不同时段的购电成本，并进一步影响储能设备的充放电安排\cite{ref4}。
因此，仅根据单一预测结果制定调度方案，可能在实际运行中出现供需失衡或购电费用增加等问题，有必要在经济性与运行风险之间进行权衡。

基于此，本文以满足小区用电需求为基本前提，综合考虑光伏出力、负荷需求、储能状态和外网电价，研究不同条件下的外网购电与储能充放电策略，并进一步分析预测偏差和电价波动对调度结果的影响。
\begin{figure}[htbp]
  \centering
  \includegraphics[width=0.7\linewidth]{background.jpg}
  \caption{并网型小区微网结构及电能流向示意图}
  \label{fig:background}
\end{figure}

\subsection{问题提出}
在上述背景下，本文以每日0:00制定当日初始购电方案为起点，按照确定性调度、实际功率变化、日内预报更新和外网电价波动四种情形，依次研究以下问题：
\begin{enumerate}
  \item \textbf{问题一：确定条件下的单日购电策略。}已知某日各时段的电价、小区负载和光伏发电预测功率，且要求储能设备在0:00与24:00的储电量相同。在此基础上，建立单日调度模型，确定各10分钟时段的计划购电量和储能充放电量，并计算全天购电量与购电费用。
  \item \textbf{问题二：考虑实际功率变化的购电策略。}已知固定的每日电价，以及每10分钟的小区负载和光伏实际功率，且紧急购电的电价是交易时刻电价的5倍。在小区负载和光伏发电功率随时间变化的条件下，制定2025年2月1日至12月31日每日0点的计划购电策略。输出上表指定日期的紧急购电结果。
  \item \textbf{问题三：考虑预报更新的购电调整策略。}根据每日0:00的光伏预报制定计划购电方案，并利用6:00、12:00和18:00的更新预报滚动调整后续购电量。综合考虑购电调整费用、违约费用和紧急购电费用，建立以总购电费用最小为目标的调度模型，给出指定日期及全年的购电策略，并分析更新预报的必要性。
  \item \textbf{问题四：波动电价下的购电策略.}针对外网电价实时波动的情况，重新建立微网购电调度模型，并分别求解问题二和问题三，分析波动电价对购电策略及购电费用的影响，给出相应结果并保存完整购电方案。
\end{enumerate}
